%------------------------------------------------------------------------------%
\documentclass[fleqn,utf8,aspectratio=169,12pt]{beamer}

\usepackage{integracao_e_series}

\title{Séries de Taylor}

%------------------------------------------------------------------------------%
\begin{document}

\addtitle

%------------------------------------------------------------------------------%
\section{Séries de Taylor}

%------------------------------------------------------------------------------%
\begin{frame}{Séries de Potência e Séries de Taylor}
  \begin{itemize}[<+->]
    \item \emph{Séries de Potências}                          \\[3mm]
          Dada uma Série de Potências construímos uma função  \\[11mm]
    \item \emph{Séries de Taylor}                             \\[3mm]
          Dada uma função construímos uma Série de Potências
  \end{itemize}
\end{frame}

%------------------------------------------------------------------------------%
\begin{frame}{Séries de Taylor e Séries de Maclaurin}
  \begin{itemize}[<+->]
    \item \emph{Séries de Taylor}                     \\[3mm]
    Centrada em um ponto qualquer, \emph{$a\in\R$}    \\[11mm]
    \item \emph{Séries de Maclaurin}                  \\[3mm]
    Série de Taylor centrada na origem, \emph{$a=0$}
  \end{itemize}
\end{frame}

%------------------------------------------------------------------------------%
\section{Coeficientes da Série de Taylor}

%------------------------------------------------------------------------------%
\begin{frame}{Construção da Série de Taylor}
  Supomos que uma função pode ser escrita como uma Série de Potências
  \[
    f(x) = \sum_{n=0}^{\infty} c_n (x-a)^n
  \]
\end{frame}

%------------------------------------------------------------------------------%
\begin{frame}[t]{Construção da Série de Taylor}
  \vspace{\baselineskip}
  \begin{align*}
    \onslide<+->{f(x)                 & = \sum_{n=0}^{\infty} c_n (x-a)^n      \\[2mm]}
    \onslide<+->{\phantom{f^{(3)}(x)} & = \memph{c_0} + c_1(x-a) + c_2(x-a)^2 + \cdots \\[5mm]}
    \onslide<+->{f(a)                 & = \memph{c_0}}
  \end{align*}
  \onslide<+->{
  \[ \memph{c_0} = f(a) = \frac{f(a)}{0!} \]}
\end{frame}

%------------------------------------------------------------------------------%
\begin{frame}[t]{Construção da Série de Taylor}
  \vspace{\baselineskip}
  \begin{align*}
    \onslide<+->{f'(x)
      & \vphantom{\sum_{n=0}^{\infty}}
        = \frac{d}{dx}\Big[c_0
        + \memph{c_1}(x-a)
        + c_2(x-a)^2
        + \cdots\Big] \\[2mm]}
    \onslide<+->{\phantom{f^{(3)}(x)}
      & = \memph{c_1}
        + 2c_2(x-a)
        + 3c_3(x-a)^2
        + \cdots \\[5mm]}
    \onslide<+->{f'(a) & = \memph{c_1}}
  \end{align*}
  \onslide<+->{
  \[
    \memph{c_1} = f'(a) = \frac{f'(a)}{1!}
  \]}
\end{frame}

%------------------------------------------------------------------------------%
\begin{frame}[t]{Construção da Série de Taylor}
  \vspace{\baselineskip}
  \begin{align*}
    \onslide<+->{f''(x)
      & \vphantom{\sum_{n=0}^{\infty}}
        = \frac{d}{dx}\Big[ c_1
        + 2\memph{c_2}(x-a)
        + 3c_3(x-a)^2
        + \cdots\Big] \\[2mm]}
    \onslide<+->{\phantom{f^{(3)}(x)}
      & = (2\times 1)\memph{c_2}
        + (3\times 2)c_3(x-a)
        + (4\times 3)c_4(x-a)^2
        + \cdots \\[5mm]}
    \onslide<+->{f''(a)
      & = (2\times 1)\memph{c_2}}
  \end{align*}
  \onslide<+->{\[
    \memph{c_2} = \frac{f''(a)}{2!}
  \]}
\end{frame}

%------------------------------------------------------------------------------%
\begin{frame}[t]{Construção da Série de Taylor}
  \vspace{\baselineskip}
  \begin{align*}
    \onslide<+->{f^{(3)}(x)
      & \vphantom{\sum_{n=0}^{\infty}}
        = \frac{d}{dx}\Big[(2\times 1)c_2
        + (3\times 2)\memph{c_3}(x-a)
        + (4\times 3)c_4(x-a)^2
        + \cdots\Big]
      \\[2mm]}
    \onslide<+->{\phantom{f^{(3)}(x)}
      & = (3\times 2\times 1)\memph{c_3}
        + (4\times 3\times 2)c_4(x-a)
        + (5\times 4\times 3)c_5(x-a)^2
        + \cdots \\[5mm]}
    \onslide<+->{f^{(3)}(a)
      & = (3\times 2\times 1)\memph{c_3}}
  \end{align*}
  \onslide<+->{
  \[
    \memph{c_3} = \frac{f^{(3)}(a)}{3!}
  \]}
\end{frame}

%------------------------------------------------------------------------------%
\begin{frame}[t]{Construção da Série de Taylor}
  \vspace{\baselineskip}
  \begin{align*}
    \onslide<+->{f^{(4)}(x)
      &  \vphantom{\sum_{n=0}^{\infty}}
        = \frac{d}{dx}\Big[(3\times 2\times 1)c_3
        + (4\times 3\times 2)\memph{c_4}(x-a)
        + (5\times 4\times 3)c_5(x-a)^2
        + \cdots\Big]
       \\[2mm]}
    \onslide<+->{\phantom{f^{(3)}(x)}
      & = (4\times 3\times 2\times 1)\memph{c_4}
        + (5\times 4\times 3\times 2)c_5(x-a)
        + \cdots \\[5mm]}
    \onslide<+->{f^{(4)}(a)
      & = (4\times 3\times 2\times 1)\memph{c_4}}
  \end{align*}
  \onslide<+->{
  \[
    \memph{c_4} = \frac{f^{(4)}(a)}{4!}
  \]}
\end{frame}

%------------------------------------------------------------------------------%
\begin{frame}{Construção da Série de Taylor}
  Generalizando
  \[
    f^{(n)}(a) = n!c_n  \qquad\qquad c_n = \frac{\;f^{(n)}(a)\;}{n!}
  \]

  \pause
  \vspace{\baselineskip}
  Se    \(\quad f(x) = \sum_{n=0}^{\infty} c_n (x-a)^n \quad \)
  então \(\quad c_n = \frac{\;f^{(n)}(a)\;}{n!} \quad \)
\end{frame}

%------------------------------------------------------------------------------%
\section{Definição da Série de Taylor}

%------------------------------------------------------------------------------%
\begin{frame}{Série de Taylor}
  Se $f$ é uma função infinitamente derivável em \(x=a\) \\
  sua \emph{Série de Taylor} centrada em $\;x=a\;$ é
  \begin{multline*}
    \sum_{n=0}^\infty \frac{\;f^{(n)\;}(a)}{n!}\,(x-a)^n = \\
    f(a)\;+\;f'(a)(x-a)\;+\;\frac{f''(a)}{2!}(x-a)^2\;+\cdots+\;\frac{f^{(n)}(a)}{n!}(x-a)^n\;+\cdots
  \end{multline*}

  \pause
  \vspace{\baselineskip}
  Observe que \emph{não} escrevemos que $f(x)$ é \emph{igual} a sua série de Taylor
\end{frame}

%------------------------------------------------------------------------------%
\begin{frame}{Série de Maclaurin}
  Uma Série de Maclaurin é uma Série de Taylor centrada em zero, $a=0$
  \[
    \sum_{n=0}^\infty \frac{\;f^{(n)\;}(0)}{n!}\,x^n
  \]
\end{frame}

%------------------------------------------------------------------------------%
\begin{frame}{Polinômios de Taylor}
  Se $f$ é uma função com $k$ derivadas em \(x=a\) \\
  seu \emph{Polinômio de Taylor} de ordem $k$ centrado em $\;x=a\;$ é
  \[
    \sum_{n=0}^k \frac{\;f^{(n)\;}(a)}{n!}\,(x-a)^n
      \;=\; f(a)
      \;+\;f'(a)(x-a)
      \;+\cdots+\; \frac{f^{(k)}(a)}{k!}(x-a)^k
    \]
\end{frame}

%------------------------------------------------------------------------------%
\section{Exemplos}

%------------------------------------------------------------------------------%
\addtocounter{exemplo}{1}
\begin{frame}{Exemplo \theexemplo}
  Calcule a Série de Taylor da função \emph{$f(x)=e^x$} centrada em \emph{$a=0$}
\end{frame}

%------------------------------------------------------------------------------%
\begin{frame}{Exemplo \theexemplo\ -- Derivadas}
  Derivando
  \[
    f^{(k)}(x) = \frac{d^k}{dx^k}\,e^x = e^x
  \]
  \pause
  Avaliando em $a=0$
  \[
    f^{(k)}(0) = e^0 = 1
  \]
\end{frame}

%------------------------------------------------------------------------------%
\begin{frame}{Exemplo \theexemplo\ -- Calculando os Coeficientes}
  \begin{align*}
    \onslide<+->{c_0 & = \frac{f^{(0)}(0)}{0!} = \frac{1}{0!} = 1  \\[3mm]}
    \onslide<+->{c_1 & = \frac{f^{(1)}(0)}{1!} = \frac{1}{1!} = 1 \\[3mm]}
    \onslide<+->{c_2 & = \frac{f^{(2)}(0)}{2!} = \frac{1}{2!} = \frac{1}{2} \\[7mm]}
    \onslide<+->{c_k & = \frac{f^{(k)}(0)}{k!} = \frac{1}{k!}}
  \end{align*}
\end{frame}

%------------------------------------------------------------------------------%
\begin{frame}{Exemplo \theexemplo\ -- Construindo a Série}
  \begin{align*}
    \onslide<+->{S(x) & = \sum_{n=0}^\infty \frac{f^{(n)}(a)}{n!}(x-a)^n \\[3mm]}
    \onslide<+->{     & = \sum_{n=0}^\infty \frac{f^{(n)}(0)}{n!}x^n     \\[3mm]}
    \onslide<+->{     & = \sum_{n=0}^\infty \frac{x^n}{n!}               \\[3mm]}
    \onslide<+->{     & = 1 + x + \frac{x^2}{2} + \frac{x^3}{3!} + \cdots + \frac{x^n}{n!}+\cdots}
  \end{align*}
\end{frame}

%------------------------------------------------------------------------------%
\begin{frame}{Exemplo \theexemplo\ -- Convergência}
\begin{itemize}[<+->]
  \item Vimos que a Série de Potências converge para todo $x\in\R$ \\[5mm]
  \item Mas a série converge para a função? \\[5mm]
  \[
     e^x \;\;{\color{red}\stackrel{?}{ = }}\;\; \sum_{n=0}^\infty \frac{\;x^n\;}{n!}
  \]
\end{itemize}
\end{frame}

%------------------------------------------------------------------------------%
\begin{frame}{Polinômios de Taylor a Função Exponencial}
  \vspace{\baselineskip}
  \scalebox{0.85}{% Title: gl2ps_renderer figure
% Creator: GL2PS 1.4.2, (C) 1999-2020 C. Geuzaine
% For: Octave
% CreationDate: Tue Sep 28 14:08:54 2021
\setlength{\unitlength}{1pt}
\begin{picture}(0,0)
\includegraphics[width=17cm,height=7cm]{fig_taylor_exp}
\end{picture}%
\begin{picture}(481,198)(0,0)
\fontsize{10}{0}\selectfont\put(21.9528,28.661){\makebox(0,0)[t]{\textcolor[rgb]{0.15,0.15,0.15}{{-3}}}}
\fontsize{10}{0}\selectfont\put(107.622,28.661){\makebox(0,0)[t]{\textcolor[rgb]{0.15,0.15,0.15}{{-2}}}}
\fontsize{10}{0}\selectfont\put(193.291,28.661){\makebox(0,0)[t]{\textcolor[rgb]{0.15,0.15,0.15}{{-1}}}}
\fontsize{10}{0}\selectfont\put(364.63,28.661){\makebox(0,0)[t]{\textcolor[rgb]{0.15,0.15,0.15}{{1}}}}
\fontsize{10}{0}\selectfont\put(450.299,28.661){\makebox(0,0)[t]{\textcolor[rgb]{0.15,0.15,0.15}{{2}}}}
\fontsize{10}{0}\selectfont\put(283.956,74.207){\makebox(0,0)[l]{\textcolor[rgb]{0.15,0.15,0.15}{{1}}}}
\fontsize{10}{0}\selectfont\put(283.956,112.272){\makebox(0,0)[l]{\textcolor[rgb]{0.15,0.15,0.15}{{2}}}}
\fontsize{10}{0}\selectfont\put(283.956,150.337){\makebox(0,0)[l]{\textcolor[rgb]{0.15,0.15,0.15}{{3}}}}
\fontsize{10}{0}\selectfont\put(283.956,188.403){\makebox(0,0)[l]{\textcolor[rgb]{0.15,0.15,0.15}{{4}}}}
\fontsize{10}{0}\selectfont\put(119.616,9.49606){\makebox(0,0)[l]{\textcolor[rgb]{0,0,0}{{$P_1$}}}}
\fontsize{10}{0}\selectfont\put(66.5008,9.49606){\makebox(0,0)[l]{\textcolor[rgb]{0,0,0}{{$P_3$}}}}
\fontsize{10}{0}\selectfont\put(3.9622,9.49606){\makebox(0,0)[l]{\textcolor[rgb]{0,0,0}{{$P_5$}}}}
\fontsize{10}{0}\selectfont\put(11.6725,98.9494){\makebox(0,0)[r]{\textcolor[rgb]{0,0,0}{{$P_2$}}}}
\fontsize{10}{0}\selectfont\put(11.6725,140.821){\makebox(0,0)[r]{\textcolor[rgb]{0,0,0}{{$P_4$}}}}
\end{picture}
}
\end{frame}

%------------------------------------------------------------------------------%
\addtocounter{exemplo}{1}
\begin{frame}{Exemplo \theexemplo}
  Calcule a série de Taylor do logaritmo natural, \(\ln(x)\), centrada em 1
\end{frame}

%------------------------------------------------------------------------------%
\begin{frame}{Exemplo \theexemplo -- Derivadas}
  \onslide<+->{Derivando}
  \begin{align*}
    \onslide<+->{f^{(0)}(x) & = \ln(x) \\}
    \onslide<+->{f^{(1)}(x) & = x^{-1} \\}
    \onslide<+->{f^{(2)}(x) & = (-1) x^{-2} \\}
    \onslide<+->{f^{(3)}(x) & = (-1)(-2) x^{-3} \\}
    \onslide<+->{f^{(4)}(x) & = (-1)(-2)(-3) x^{-4} \\}
    \onslide<+->{f^{(5)}(x) & = (-1)(-2)(-3)(-4) x^{-5} \\[7mm]}
    \onslide<+->{f^{(n)}(x) & = (-1)^{n-1} \, (n-1)!\, x^{-n} \qquad \qquad n = 1, 2, 3, \dots}
  \end{align*}
\end{frame}

%------------------------------------------------------------------------------%
\begin{frame}{Exemplo \theexemplo\ -- Derivadas}
  Avaliando as derivadas em $x=1$

  \pause
  Para $n=0$
  \[
    \pause
    f^{(0)}(1)
    \pause
    = \ln(1)
    \pause
    = 0
  \]

  \pause
  Para $n = 1, 2, 3, \dots$
  \[
    \pause
    f^{(n)}(1)
    \pause
    = (-1)^{n-1} \, (n-1)!\, 1^{-n}
    \pause
    = (-1)^{n-1} \, (n-1)!
  \]
\end{frame}

%------------------------------------------------------------------------------%
\begin{frame}{Exemplo \theexemplo\ -- Calculando os Coeficientes}
  Calculando os coeficientes $c_n$

  \pause
  Para $n=0$
  \[
    \pause
    c_0
    \pause
    = \frac{f^{(0)}(1)}{0!}
    \pause
    = \frac{\ln(1)}{1}
    \pause
    = \frac{0}{1}
    \pause
    = 0
  \]

  \pause
  Para $n = 1, 2, 3, \dots$
  \[
    \pause
    c_n
    \pause
    = \frac{f^{(n)}(1)}{n!}
    \pause
    = \frac{(-1)^{n-1} \, (n-1)!}{n!}
    \pause
    = \frac{(-1)^{n-1}}{n}
    \pause
    = \frac{(-1)^{n+1}}{n}
  \]
\end{frame}

%------------------------------------------------------------------------------%
\begin{frame}{Exemplo \theexemplo\ -- Construindo a Série}
  \begin{align*}
    \onslide<+->{S(x) & = \sum_{n=0}^\infty \frac{f^{(n)}(a)}{n!}(x-a)^n \\[1.5mm]}
    \onslide<+->{& = \sum_{n=0}^\infty \frac{f^{(n)}(1)}{n!}(x-1)^n      \\[1.5mm]}
    \onslide<+->{& = \sum_{n=1}^\infty \frac{(-1)^{n+1}}{n} (x-1)^n      \\[1.5mm]}
    \onslide<+->{& = \sum_{n=1}^\infty (-1)^{n+1} \frac{(x-1)^n}{n}      \\[1.5mm]}
    \onslide<+->{& = (x-1)
                   - \frac{(x-1)^2}{2}
                   + \frac{(x-1)^3}{3}
                   + \cdots
                   + (-1)^{n+1}\frac{(x-1)^n}{n!}
                   + \cdots}
  \end{align*}
\end{frame}

%------------------------------------------------------------------------------%
\addtocounter{exemplo}{1}
\begin{frame}{Exemplo \theexemplo}
  Calcule a série de Taylor do cosseno, \(\cos(x)\), centrada em 0
\end{frame}

%------------------------------------------------------------------------------%
\begin{frame}{Exemplo \theexemplo -- Derivadas}
\begin{columns}
  \column{0.5\linewidth}
  \onslide<+->{Derivando}
  \begin{align*}
    \onslide<+->{f^{(0)}(x) & = \hpm \cos(x) \\}
    \onslide<+->{f^{(1)}(x) & =     -\sin(x) \\}
    \onslide<+->{f^{(2)}(x) & =     -\cos(x) \\}
    \onslide<+->{f^{(3)}(x) & = \hpm \sin(x) \\}
    \onslide<+->{f^{(4)}(x) & = \hpm \cos(x) \\}
    \onslide<+->{f^{(5)}(x) & =     -\sin(x) \\}
    \onslide<+->{f^{(6)}(x) & =     -\cos(x)}
  \end{align*}
  \column{0.5\linewidth}
  \onslide<+->{Avaliando as derivadas em $x=0$}
  \begin{align*}
    \onslide<+->{f^{(0)}(0) & = \hpm \cos(0) = \hpm 1  \\}
    \onslide<+->{f^{(1)}(0) & =     -\sin(0) = \hpm 0  \\}
    \onslide<+->{f^{(2)}(0) & =     -\cos(0) =     -1  \\}
    \onslide<+->{f^{(3)}(0) & = \hpm \sin(0) = \hpm 0  \\}
    \onslide<+->{f^{(4)}(0) & = \hpm \cos(0) = \hpm 1  \\}
    \onslide<+->{f^{(5)}(0) & =     -\sin(0) = \hpm 0  \\}
    \onslide<+->{f^{(6)}(0) & =     -\cos(0) =      -1   }
  \end{align*}
\end{columns}
\end{frame}

%------------------------------------------------------------------------------%
\begin{frame}{Exemplo \theexemplo\ -- Representando o padrão}
\begin{columns}
  \column{0.5\linewidth}
  Para $n$ ímpar
  \[
    f^{(n)}(0) = 0
  \]

  \pause
  Para $n$ par
  \begin{align*}
    f^{(0)}(0) = \hpm 1  \\
    f^{(2)}(0) =     -1  \\
    f^{(4)}(0) = \hpm 1  \\
    f^{(6)}(0) =     -1
  \end{align*}

  \column{0.5\linewidth}
  \pause
  Para todo $k$, \(n=2k\) é sempre par
  \[
    \begin{array}{ccc}
      \hline
      \hpm k \hpm & \hpm n \hpm & f^{(2k)}(0) \\ \hline
      0           & 0           & \hpm 1 \hpm \\
      1           & 2           &     -1 \hpm \\
      2           & 4           & \hpm 1 \hpm \\
      3           & 6           &     -1 \hpm \\
      4           & 8           & \hpm 1 \hpm \\ \hline
    \end{array}
  \]

  \pause
  Assim
  \(
    f^{(2k)}(0) = (-1)^k
  \)
\end{columns}
\end{frame}

%------------------------------------------------------------------------------%
\begin{frame}{Exemplo \theexemplo\ -- Calculando os Coeficientes}
  Calculando os coeficientes $c_n$

  \pause
  Para $n$ ímpar
  \[
    c_n = 0
  \]

  \pause
  Para $n$ par, usamos $n=2k$
  \[
    \pause
    c_n
    \pause
    = c_{2k}
    \pause
    = \frac{f^{(n)}(0)}{n!}
    \pause
    = \frac{f^{(2k)}(0)}{(2k)!}
    \pause
    = \frac{(-1)^k}{(2k)!}
  \]
\end{frame}

%------------------------------------------------------------------------------%
\begin{frame}{Exemplo \theexemplo\ -- Construindo a Série}
  \begin{align*}
    \onslide<+->{S(x) & = \sum_{n=0}^\infty \frac{f^{(n)}(a)}{n!}(x-a)^n \\[1.5mm]}
    \onslide<+->{& = \sum_{n=0}^\infty \frac{f^{(n)}(0)}{n!}x^n          \\[1.5mm]}
    \onslide<+->{& = \sum_{k=0}^\infty \frac{f^{(2k)}(0)}{(2k)!}x^{2k}   \\[1.5mm]}
    \onslide<+->{& = \sum_{k=0}^\infty (-1)^k\frac{x^{2k}}{(2k)!}        \\[1.5mm]}
    \onslide<+->{& = 1
                   - \frac{x^2}{2}
                   + \frac{x^4}{4!}
                   - \frac{x^6}{6!}
                   + \cdots
                   + (-1)^k \frac{x^{2k}}{(2k)!}
                   + \cdots}
  \end{align*}
\end{frame}

%------------------------------------------------------------------------------%
\addtocounter{exemplo}{1}
\begin{frame}{Exemplo \theexemplo}
  Encontre a série de Maclaurin para
  \(f(x) = \frac{1}{(1-x)^2}\)
  e o raio de convergência da série de potências.
\end{frame}

%------------------------------------------------------------------------------%
\begin{frame}{Exemplo \theexemplo\ -- Calculando as derivadas}
\begin{columns}
\column{0.5\linewidth}
\begin{align*}
  \onslide<+->{f^{(0)}(x) & = (1-x)^{-2}                      \\}
  \onslide<+->{f^{(1)}(x) & = (1-x)^{-3} \,2                  \\}
  \onslide<+->{f^{(2)}(x) & = (1-x)^{-4} \,2 \cdot 3          \\}
  \onslide<+->{f^{(3)}(x) & = (1-x)^{-5} \,2 \cdot 3 \cdot 4  \\[7mm]}
  \onslide<+->{f^{(n)}(x) & = (1-x)^{-n-2} (n+1)!}
\end{align*}
\column{0.5\linewidth}
\begin{align*}
  \onslide<+->{f^{(0)}(0) & = 1      \\}
  \onslide<+->{f^{(1)}(0) & = 2     \\}
  \onslide<+->{f^{(2)}(0) & = 3!     \\}
  \onslide<+->{f^{(3)}(0) & = 4!     \\[7mm]}
  \onslide<+->{f^{(n)}(0) & = (n+1)!}
\end{align*}
\end{columns}
\end{frame}

%------------------------------------------------------------------------------%
\begin{frame}{Exemplo \theexemplo\ -- Série de Maclaurin}
  \[
    \sum_{n=0}^\infty \frac{f^{(n)}(0)}{n!} x^n
    \pause \;=\; \sum_{n=0}^\infty \frac{(n+1)! }{n!}x^n
    \pause \;=\; \sum_{n=0}^\infty \frac{(n+1)\,n! }{n!}x^n
    \pause \;=\; \sum_{n=0}^\infty (n+1) x^n
  \]
\end{frame}

%%------------------------------------------------------------------------------%
%\begin{frame}{Exemplo \theexemplo -- Analisando a convergência}
%  Usando o teste da razão para determinar o Raio de Convergência
%  \pause
%  \[
%    a_n = (n+1) x^n
%    \hspace{57.2mm}
%    \pause
%    \abs{a_n} = (n+1) \abs{x}^n
%  \]
%  \[
%    \pause
%    a_{n+1}
%    \pause
%    = [(n+1)+1] x^{n+1}
%    \pause
%    = (n+2) x^{n+1}
%    \hspace{15mm}
%    \pause
%    \abs{a_{n+1}} = (n+2) \abs{x}^{n+1}
%  \]
%
%  \pause
%  \vspace{0.1\baselineskip}
%  \[
%    \rho
%    = \lim_{n\to\infty}\frac{\abs{a_{n+1}}}{\abs{a_n}}
%    \pause
%    = \lim_{n\to\infty}\frac{(n+2)\abs{x}^{n+1}}{(n+1)\abs{x}^n}
%    \pause
%    = \abs{x}\lim_{n\to\infty}\frac{n+2}{n+1}
%    \pause
%    = \abs{x}\lim_{n\to\infty}\frac{1+\sfrac{2}{n}}{1+\sfrac{1}{n}}
%    \pause
%    = \abs{x}
%  \]
%
%  \pause
%  \vspace{0.5\baselineskip}
%  A série converge para
%  \(\rho = \abs{x} < 1\)
%  \pause
%  portanto \(R=1\)
%\end{frame}


%------------------------------------------------------------------------------%
\begin{frame}{Exemplo \theexemplo\ -- Analisando a convergência}
  Usando o teste da razão para determinar o Raio de Convergência
  \begin{align*}
      \onslide<+->{a_n & = (n+1) x^n & }
      \onslide<+->{\abs{a_n} & = (n+1) \abs{x}^n}
      \\[3mm]
      \onslide<+->{a_{n+1}}
      \onslide<+->{& = [(n+1)+1] x^{n+1}}
      \onslide<+->{= (n+2) x^{n+1}}
      \onslide<+->{& \abs{a_{n+1}} & = (n+2) \abs{x}^{n+1}}
  \end{align*}

  \[
    \onslide<+->{\rho}
    \onslide<+->{= \lim_{n\to\infty}\frac{\abs{a_{n+1}}}{\abs{a_n}}}
    \onslide<+->{= \lim_{n\to\infty}\frac{(n+2)\abs{x}^{n+1}}{(n+1)\abs{x}^n}}
    \onslide<+->{= \abs{x}\lim_{n\to\infty}\frac{n+2}{n+1}}
    \onslide<+->{= \abs{x}\lim_{n\to\infty}\frac{1+\sfrac{2}{n}}{1+\sfrac{1}{n}}}
    \onslide<+->{= \abs{x}}
  \]

  \vspace{0.5\baselineskip}
  \onslide<+->{A série converge para \(\rho = \abs{x} < 1\)}
  \onslide<+->{portanto \(R=1\)}
\end{frame}


%------------------------------------------------------------------------------%
\section{Lista Mínima}

%------------------------------------------------------------------------------%
\begin{frame}{Lista Mínima}
  Estudar a Seção 8.1 da Apostila

  \vfill
  Exercícios: 1a-d, 2d-h, 3a-d, 5a-c

  \vfill
  Atenção: A prova é baseada no livro, não nas apresentações
\end{frame}

%------------------------------------------------------------------------------%
\end{document}
%------------------------------------------------------------------------------%
